% titel: Blatt 9, A3: Rationale Funktionen
% Teil II: der durchgerechnete Weg.

\subsection*{a) Stammfunktionen per Partialbruch}

\textbf{i) $f(x)=\dfrac{x+1}{x(x-5)}$:} Ansatz
$\frac{x+1}{x(x-5)}=\frac{A}{x}+\frac{B}{x-5}$. Mit $x(x-5)$ multiplizieren:
\begin{equation}
  x+1 = A(x-5)+Bx .
\end{equation}
$x=0$: $1=-5A\Rightarrow A=-\frac15$. $x=5$: $6=5B\Rightarrow B=\frac65$.
\begin{align}
  \int \frac{x+1}{x(x-5)}\,\dd x
  &= -\frac15\int\frac1x\,\dd x + \frac65\int\frac{1}{x-5}\,\dd x \\
  &= -\frac15\ln\abs{x} + \frac65\ln\abs{x-5} + C .
\end{align}

\textbf{ii) $f(x)=\dfrac{x^2-2x}{(x-1)^2(x+3)}$:} doppelte Nullstelle bei $1$:
\begin{equation}
  \frac{x^2-2x}{(x-1)^2(x+3)}
  = \frac{A}{x-1} + \frac{B}{(x-1)^2} + \frac{C}{x+3}.
\end{equation}
Mit $(x-1)^2(x+3)$ multiplizieren:
\begin{equation}
  x^2-2x = A(x-1)(x+3) + B(x+3) + C(x-1)^2 .
\end{equation}
$x=1$: $-1 = 4B \Rightarrow B=-\frac14$. $x=-3$: $15=16C\Rightarrow
C=\frac{15}{16}$. $x=0$: $0=-3A+3B+C$
$\Rightarrow 3A=3B+C=-\frac34+\frac{15}{16}=-\frac{12}{16}+\frac{15}{16}=\frac{3}{16}$
$\Rightarrow A=\frac{1}{16}$.
\begin{align}
  \int \frac{x^2-2x}{(x-1)^2(x+3)}\,\dd x
  &= \frac1{16}\int\frac1{x-1}\,\dd x
   - \frac14\int\frac1{(x-1)^2}\,\dd x
   + \frac{15}{16}\int\frac1{x+3}\,\dd x \\
  &= \frac1{16}\ln\abs{x-1}
   + \frac{1}{4(x-1)}
   + \frac{15}{16}\ln\abs{x+3} + C .
\end{align}

\begin{ergebnis}
\[
\int\frac{x+1}{x(x-5)}\,\dd x=-\frac15\ln\abs{x}+\frac65\ln\abs{x-5}+C,
\qquad
\int\frac{x^2-2x}{(x-1)^2(x+3)}\,\dd x
=\frac1{16}\ln\abs{x-1}+\frac{1}{4(x-1)}+\frac{15}{16}\ln\abs{x+3}+C.
\]
\end{ergebnis}

\subsection*{b) Uneigentliches Integral}
Substituiere $u=2x-1$, $\dd u=2\dd x$; $x=1\Rightarrow u=1$,
$x\to\infty\Rightarrow u\to\infty$:
\begin{align}
  \int_1^\infty \frac{\dd x}{(2x-1)^4}
  &= \frac12\int_1^\infty u^{-4}\,\dd u
   = \frac12\Bigl[-\frac13 u^{-3}\Bigr]_1^\infty
   = -\frac16\Bigl(\lim_{R\to\infty}\frac1{R^3}-1\Bigr)
   = \frac16 .
\end{align}

\begin{ergebnis}
$\displaystyle\int_1^\infty\frac{\dd x}{(2x-1)^4}=\frac16 .$
\end{ergebnis}

\begin{kontrolle}
\textbf{Plausibilität:} alle Integranden fallen wie $x^{-2}$ bzw.\ $x^{-4}$
ab; das uneigentliche Integral ist endlich und klein. Durch Differenzieren
der gefundenen Stammfunktionen erhält man wieder $f$ (Koeffizientenprobe).
\end{kontrolle}

\begin{merksatz}
Partialbruchzerlegung: einfache Nullstelle $\to\frac{A}{x-a}$, $m$-fache
$\to$ Summe bis $\frac{A_m}{(x-a)^m}$; Koeffizienten über geschickte
$x$-Werte. Uneigentliches Integral: substituieren, Stammfunktion, dann
Grenzwert.
\end{merksatz}
